\begin{summarybox}
\textbf{\textcolor{CSummary}{一句话核心}}

\textbf{连不等式 $N < f(x) < M$ 可等价转化为乘积、绝对值、分式、倒数四种形式，本质是函数值落在开区间内。}

\medskip
\textbf{\textcolor{CSummary}{使用条件}}
\begin{itemize}[leftmargin=*, itemsep=0.5em, topsep=0.4em]
    \item \textbf{条件 1（实数大小）：} $M$ 和 $N$ 为实数，且满足 $M > N$。
    \item \textbf{条件 2（函数实值）：} 函数 $f(x)$ 在相关定义域上取实数值。
\end{itemize}

\medskip
\textbf{\textcolor{CSummary}{关键提醒}}
\begin{itemize}[leftmargin=*, itemsep=0.5em, topsep=0.4em]
    \item \textbf{易错点（倒数条件）：} 使用倒数形式时，先确认分母非零及正值；在 $M>N$ 下，倒数不等式本身可推出正值。
    \item \textbf{检查项（形式适用）：} 在证明题中优先考虑乘积或绝对值形式；在求解题中可尝试分式或倒数形式简化计算。
\end{itemize}
\medskip
\textbf{方法选择表（$t=f(x)$）}
\begin{center}
\begin{tabular}{lll}
形式 & 识别特征 & 适用场景\\\hline
原始 & 上下界 & 直接移项\\
乘积 & 两个端点差 & 因式分解、符号判断\\
绝对值 & 中点与半长 & 距离、配方\\
分式 & 端点差之比 & 分式符号\\
倒数 & 两个正数的倒数 & 倒数比较
\end{tabular}
\end{center}
\textbf{记忆：} 夹在中间，乘积为负；寻找中点，转化距离；分子分母，符号相同；正数倒数，大小反转。

选择变形的目的不是增加计算，而是利用已有结构缩短解题路径。
\end{summarybox}
